\documentclass[a4paper,12pt,twoside]{memoir}

% Castellano
\usepackage[spanish,es-tabla]{babel}
\selectlanguage{spanish}
\usepackage[utf8]{inputenc}
\usepackage[T1]{fontenc}
\usepackage{lmodern} % scalable font
\usepackage{microtype}
\usepackage{placeins}

\RequirePackage{booktabs}
\RequirePackage[table]{xcolor}
\RequirePackage{xtab}
\RequirePackage{multirow}

% Links
\usepackage[colorlinks]{hyperref}
\hypersetup{
	allcolors = {red}
}

% Ecuaciones
\usepackage{amsmath}

% Rutas de fichero / paquete
\newcommand{\ruta}[1]{{\sffamily #1}}

% Párrafos
\nonzeroparskip


% Imagenes
\usepackage{graphicx}
\newcommand{\imagen}[2]{
	\begin{figure}[!h]
		\centering
		\includegraphics[width=0.9\textwidth]{#1}
		\caption{#2}\label{fig:#1}
	\end{figure}
	\FloatBarrier
}

\newcommand{\imagenflotante}[2]{
	\begin{figure}%[!h]
		\centering
		\includegraphics[width=0.9\textwidth]{#1}
		\caption{#2}\label{fig:#1}
	\end{figure}
}



% El comando \figura nos permite insertar figuras comodamente, y utilizando
% siempre el mismo formato. Los parametros son:
% 1 -> Porcentaje del ancho de página que ocupará la figura (de 0 a 1)
% 2 --> Fichero de la imagen
% 3 --> Texto a pie de imagen
% 4 --> Etiqueta (label) para referencias
% 5 --> Opciones que queramos pasarle al \includegraphics
% 6 --> Opciones de posicionamiento a pasarle a \begin{figure}
\newcommand{\figuraConPosicion}[6]{%
  \setlength{\anchoFloat}{#1\textwidth}%
  \addtolength{\anchoFloat}{-4\fboxsep}%
  \setlength{\anchoFigura}{\anchoFloat}%
  \begin{figure}[#6]
    \begin{center}%
      \Ovalbox{%
        \begin{minipage}{\anchoFloat}%
          \begin{center}%
            \includegraphics[width=\anchoFigura,#5]{#2}%
            \caption{#3}%
            \label{#4}%
          \end{center}%
        \end{minipage}
      }%
    \end{center}%
  \end{figure}%
}

%
% Comando para incluir imágenes en formato apaisado (sin marco).
\newcommand{\figuraApaisadaSinMarco}[5]{%
  \begin{figure}%
    \begin{center}%
    \includegraphics[angle=90,height=#1\textheight,#5]{#2}%
    \caption{#3}%
    \label{#4}%
    \end{center}%
  \end{figure}%
}
% Para las tablas
\newcommand{\otoprule}{\midrule [\heavyrulewidth]}
%
% Nuevo comando para tablas pequeñas (menos de una página).
\newcommand{\tablaSmall}[5]{%
 \begin{table}
  \begin{center}
   \rowcolors {2}{gray!35}{}
   \begin{tabular}{#2}
    \toprule
    #4
    \otoprule
    #5
    \bottomrule
   \end{tabular}
   \caption{#1}
   \label{tabla:#3}
  \end{center}
 \end{table}
}

%
%Para el float H de tablaSmallSinColores
\usepackage{float}

%
% Nuevo comando para tablas pequeñas (menos de una página).
\newcommand{\tablaSmallSinColores}[5]{%
 \begin{table}[H]
  \begin{center}
   \begin{tabular}{#2}
    \toprule
    #4
    \otoprule
    #5
    \bottomrule
   \end{tabular}
   \caption{#1}
   \label{tabla:#3}
  \end{center}
 \end{table}
}

\newcommand{\tablaApaisadaSmall}[5]{%
\begin{landscape}
  \begin{table}
   \begin{center}
    \rowcolors {2}{gray!35}{}
    \begin{tabular}{#2}
     \toprule
     #4
     \otoprule
     #5
     \bottomrule
    \end{tabular}
    \caption{#1}
    \label{tabla:#3}
   \end{center}
  \end{table}
\end{landscape}
}

%
% Nuevo comando para tablas grandes con cabecera y filas alternas coloreadas en gris.
\newcommand{\tabla}[6]{%
  \begin{center}
    \tablefirsthead{
      \toprule
      #5
      \otoprule
    }
    \tablehead{
      \multicolumn{#3}{l}{\small\sl continúa desde la página anterior}\\
      \toprule
      #5
      \otoprule
    }
    \tabletail{
      \hline
      \multicolumn{#3}{r}{\small\sl continúa en la página siguiente}\\
    }
    \tablelasttail{
      \hline
    }
    \bottomcaption{#1}
    \rowcolors {2}{gray!35}{}
    \begin{xtabular}{#2}
      #6
      \bottomrule
    \end{xtabular}
    \label{tabla:#4}
  \end{center}
}

%
% Nuevo comando para tablas grandes con cabecera.
\newcommand{\tablaSinColores}[6]{%
  \begin{center}
    \tablefirsthead{
      \toprule
      #5
      \otoprule
    }
    \tablehead{
      \multicolumn{#3}{l}{\small\sl continúa desde la página anterior}\\
      \toprule
      #5
      \otoprule
    }
    \tabletail{
      \hline
      \multicolumn{#3}{r}{\small\sl continúa en la página siguiente}\\
    }
    \tablelasttail{
      \hline
    }
    \bottomcaption{#1}
    \begin{xtabular}{#2}
      #6
      \bottomrule
    \end{xtabular}
    \label{tabla:#4}
  \end{center}
}

%
% Nuevo comando para tablas grandes sin cabecera.
\newcommand{\tablaSinCabecera}[5]{%
  \begin{center}
    \tablefirsthead{
      \toprule
    }
    \tablehead{
      \multicolumn{#3}{l}{\small\sl continúa desde la página anterior}\\
      \hline
    }
    \tabletail{
      \hline
      \multicolumn{#3}{r}{\small\sl continúa en la página siguiente}\\
    }
    \tablelasttail{
      \hline
    }
    \bottomcaption{#1}
  \begin{xtabular}{#2}
    #5
   \bottomrule
  \end{xtabular}
  \label{tabla:#4}
  \end{center}
}



\definecolor{cgoLight}{HTML}{EEEEEE}
\definecolor{cgoExtralight}{HTML}{FFFFFF}

%
% Nuevo comando para tablas grandes sin cabecera.
\newcommand{\tablaSinCabeceraConBandas}[5]{%
  \begin{center}
    \tablefirsthead{
      \toprule
    }
    \tablehead{
      \multicolumn{#3}{l}{\small\sl continúa desde la página anterior}\\
      \hline
    }
    \tabletail{
      \hline
      \multicolumn{#3}{r}{\small\sl continúa en la página siguiente}\\
    }
    \tablelasttail{
      \hline
    }
    \bottomcaption{#1}
    \rowcolors[]{1}{cgoExtralight}{cgoLight}

  \begin{xtabular}{#2}
    #5
   \bottomrule
  \end{xtabular}
  \label{tabla:#4}
  \end{center}
}




\graphicspath{ {./img/} }

% Capítulos
\chapterstyle{bianchi}
\newcommand{\capitulo}[2]{
	\setcounter{chapter}{#1}
	\setcounter{section}{0}
	\chapter*{#2}
	\addcontentsline{toc}{chapter}{#2}
	\markboth{#2}{#2}
}

% Apéndices
\renewcommand{\appendixname}{Apéndice}
\renewcommand*\cftappendixname{\appendixname}

\newcommand{\apendice}[1]{
	%\renewcommand{\thechapter}{A}
	\chapter{#1}
}

\renewcommand*\cftappendixname{\appendixname\ }

% Formato de portada
\makeatletter
\usepackage{xcolor}
\newcommand{\tutor}[1]{\def\@tutor{#1}}
\newcommand{\course}[1]{\def\@course{#1}}
\definecolor{cpardoBox}{HTML}{E6E6FF}
\def\maketitle{
  \null
  \thispagestyle{empty}
  % Cabecera ----------------
\noindent\includegraphics[width=\textwidth]{cabecera}\vspace{1cm}%
  \vfill
  % Título proyecto y escudo informática ----------------
  \colorbox{cpardoBox}{%
    \begin{minipage}{.8\textwidth}
      \vspace{.5cm}\Large
      \begin{center}
      \textbf{TFG del Grado en Ingeniería Informática}\vspace{.6cm}\\
      \textbf{\LARGE\@title{}}
      \end{center}
      \vspace{.2cm}
    \end{minipage}

  }%
  \hfill\begin{minipage}{.20\textwidth}
    \includegraphics[width=\textwidth]{escudoInfor}
  \end{minipage}
  \vfill
  % Datos de alumno, curso y tutores ------------------
  \begin{center}%
  {%
    \noindent\LARGE
    Presentado por \@author{}\\ 
    en Universidad de Burgos --- \@date{}\\
    Tutor: \@tutor{}\\
  }%
  \end{center}%
  \null
  \cleardoublepage
  }
\makeatother


% Datos de portada
\title{Aplicación para la gestión del sonido en un dispositivo móvil \\Documentación Técnica}
\author{Jorge Zamora Marqués}
\tutor{Carlos López Nodal}
\date{\today}

\begin{document}

\maketitle



\cleardoublepage



%%%%%%%%%%%%%%%%%%%%%%%%%%%%%%%%%%%%%%%%%%%%%%%%%%%%%%%%%%%%%%%%%%%%%%%%%%%%%%%%%%%%%%%%



\frontmatter


\clearpage

% Indices
\tableofcontents

\clearpage

\listoffigures

\clearpage

\listoftables

\clearpage

\mainmatter

\appendix

\apendice{Plan de Proyecto Software}

\section{Introducción}

\section{Planificación temporal}
En desarrollo

\section{Estudio de viabilidad}
En desarrollo

\subsection{Viabilidad económica}
En desarrollo

\subsection{Viabilidad legal}
En desarrollo



\apendice{Especificación de Requisitos}

\section{Introducción}

\section{Objetivos generales}

\section{Catalogo de requisitos}

\section{Especificación de requisitos}



\apendice{Especificación de diseño}

\section{Introducción}
En este apartado veremos cómo hemos estructurado el código y a su vez las estructuras utilizadas para el funcionamiento del proyecto.
\section{Diseño de datos}
Los datos con los que hemos estado trabajando han sido objetos, con lo cual los principales objetos estructurales con los que hemos estado trabajando han sido:
\begin{itemize}
	\item Configuración de sonido.
	\item Evento manual.
	\item Evento periódico.
	\item Evento GPS.
	\item Evento de calendario.
	\item Evento Wifi.
	\item Configuración de la app
\end{itemize}

Con estas clases se ha configurado toda la aplicación para que pueda cambiar de configuración de sonido según la especificación de requisitos del trabajo.



\section{Diseño procedimental}

Faltan las imágenes que se añadirán cuando este

\section{Diseño arquitectónico}

Como diseño arquitectónico podríamos decir que es una aplicación que se basa en el modelo vista controlador, ya que un controlador irá sacando los datos que se requieran a la vista, de un modelo de datos que no está en contacto directo con la vista, siendo el controlador el que haga el paso de la parte front-end a la back-end.



\apendice{Documentación técnica de programación}

\section{Introducción}
A continuación veremos como se ha estado trabajando para generar la aplicación y algunos conceptos necesarios.
\section{Estructura de directorios}
.\TFG-Gestor-de-sonido-para-dispositivos-Android\Codigo\app\src\main *** RUTA PRINCIPAL ***

Código java de la aplicación:
.\java
Hay que diferenciar la parte de objetos, controladores y aplicación principal.

Hay que sacar imágenes finales de la estructura.

Parte gráfica de la aplicación:
.\res
Guardamos los bitmaps de los iconos, layouts, estructuras de menús, literales de los menús...

\section{Manual del programador}
\subsection{Ficheros de configuración}
En nuestro proyecto encontraremos tres ficheros de configuración, dos ficheros que utiliza Gradle, para añadir dependencias a nuestro proyecto y lineas de ejecución de este, y un archivo manifest donde indicaremos todo lo relacionado con la ejecución de librerías de android, versionado de la aplicación y diferentes componentes android.
\subsection{Observadores}
Los observadores en nuestra aplicación son una serie de clases que se encargarán de ver si se lanza agún tipo de evento, el cual active alguna de nuestras configuraciones.

Para ello tendremos dos tipos de observadores: un observador GPS y un observador genérico.
\subsubsection{Observador GPS}
El observador GPS es la clase encargada de detectar cunado hay cambios en nuestra localización GPS, es decir, cada vez que nos movamos X metros o pase Y tiempo se lanzará un evento, el cual realizará las comprobaciones pertinentes para ver si hay configurado algún evento de GPS con esas coordenadas, y si es así activar la configuración 

\section{Compilación, instalación y ejecución del proyecto}
\section{Compilación, instalación y ejecución del proyecto}
Una vez ya realizados los pasos anteriores vistos en el manual del programador podremos realizar los siguiente pasos.
\subsection{Compilación de la aplicación}
Para ello tendremos que tener el proyecto abierto con el Android Studio y seleccionar en la barra de herramientas el botón del desplegable de "Build", y ahí seleccionaremos "Build Project" o "Rebuid Project" si ya lo habíamos buildeado anterior mente.
\subsection{Instalación de la APP y ejecución del proyecto}
Para instalar la aplicación dentro de un simulador Android de IDE Android Studio lo que haremos será crear un dispositivo virtual con la versión de Android que queramos simular. Una vez tengamos creada la máquina virtual de Android lo que haremos será pinchar en el botón de "Run" en la barra de herramientas de Android Studio, o abriendo el desplegable "Run" y seleccionando la opción "Run 'app'". Una vez hecho esto, tendremos que seleccionar un dispositivo sobre el que instalar y ejecutar la aplicación, por consiguiente se nos abrirá una ventana donde seleccionar el dispositivo simulado sobre el que correr la aplicación, permitiéndonos también crear nuevos dispositivos simulados. Una vez seleccionado el dispositivo pincharemos en el botón "OK" y esperaremos a que el proyecto se buildee, y que la máquina virtual se arranque. Una vez arrancada la máquina virtual automáticamente se instalará y se ejecutará.
\section{Pruebas del sistema}

\subsection{Pruebas unitarias de la base de datos}

\subsubsection{Creación de objetos en la base de datos}

\paragraph{Premisas}
Se tenga instanciada la base de datos y creado el modelo de datos.
Se tenga definida la clase que controle los índices de la base de datos, y esta sea arrancada cuando se inicie la aplicación.

\paragraph{Acciones}
Se generarán todos los tipos de eventos en la base de datos.
Cada vez que se genere una nueva entrada en la base de datos se comprobará que se introduce comprobando que el numero máximo del id generado es mayor al id máximo que había anteriormente.

\paragraph{Resultados esperados}

Se comprueba que se insertan correctamente todos los tipos de eventos en la base de datos.

\paragraph{Resultado}
Prueba superada.

\subsubsection{Obtención de datos de la base de datos}

\paragraph{Premisas}
Se tenga instanciada la base de datos y creado el modelo de datos.
Se tenga definida la clase que controle los índices de la base de datos, y esta sea arrancada cuando se inicie la aplicación.
Haya datos guardados dentro de la base de datos.

\paragraph{Acciones}
Se hará una búsqueda de cada tipo de datos instanciados.
Se comprobará que se puede iterar sobre dichos datos.

\paragraph{Resultados esperados}

Afirmamos que se puede iterar sobre objetos creados de la base de datos y su contenido corresponde con el introducido.

\paragraph{Resultado}
Prueba superada.

\subsection{Pruebas unitarias del observador GPS}

\subsubsection{Obtención de la localización GPS en cualquier momento}

\paragraph{Premisas}
Se tienen que tener aceptados los permisos de localización GPS en la aplicación.


\paragraph{Acciones}
Se generarán todos los tipos de eventos en la base de datos.
Cada vez que se genere una nueva entrada en la base de datos se comprobará que se introduce comprobando que el numero máximo del id generado es mayor al id máximo que había anteriormente.

\paragraph{Resultados esperados}

Se comprueba que se insertan correctamente todos los tipos de eventos en la base de datos.

\paragraph{Resultado}
Prueba superada.

\subsubsection{Acceder a los eventos de calendario}

\paragraph{Premisas}
Se tendrán los permisos de lectura y escritura de calendarios en la aplicación.


\paragraph{Acciones}
Se obtendrán los calendarios y sus eventos para iterar sobre ellos.

\paragraph{Resultados esperados}
Se esperará un listado de eventos sobre el que poder iterar.

\paragraph{Resultado}
Prueba superada.


\subsubsection{Obtener SSID del wifi}

\paragraph{Premisas}
No se requieren premisas.

\paragraph{Acciones}
Se obtendrá la ssid del wifi actual.

\paragraph{Resultados esperados}
Se obtendrá unknown ssid ya que no se puede conectar el emulador a redes wifi.

\paragraph{Resultado}
Prueba superada.

\subsubsection{Observador genérico}

\paragraph{Premisas}
No se requieren premisas.

\paragraph{Acciones}
Se tendrán que crear eventos para ver si la aplicación los detecta.

\paragraph{Resultados esperados}
Si se encuentran eventos que coincidan con las condiciones, se tendrá que informar de ellos.

\paragraph{Resultado}
Prueba superada.

\subsection{Pruebas automáticas de la aplicación}



\apendice{Documentación de usuario}

\section{Introducción}

\section{Requisitos de usuarios}

\section{Instalación}

\section{Manual del usuario}





\bibliographystyle{plain}
\bibliography{bibliografiaAnexos}

\end{document}
